\documentclass[11pt]{article}
\usepackage[margin=1in]{geometry}
% \usepackage{microtype}
\usepackage{listings}
\usepackage{xcolor}
\usepackage{hyperref}
\usepackage{enumitem}
\hypersetup{colorlinks=true, linkcolor=blue!50!black, urlcolor=blue!60!black}

\lstdefinestyle{code}{
  basicstyle=\small\ttfamily,
  keywordstyle=\color{blue!60!black},
  commentstyle=\color{gray},
  stringstyle=\color{red!50!black},
  showstringspaces=false,
  breaklines=true,
  numbers=left, numbersep=5pt, numberstyle=\tiny\color{gray},
  frame=single, framesep=4pt,
}
\lstset{style=code}

\title{A Critical Evaluation of \texttt{dbdreader} \\
       \large and a Pure-Julia Replacement (\texttt{JLDBDReader.jl})}
\author{Translation analysis and reimplementation notes}
\date{}

\begin{document}
\maketitle

\section{Scope}

The Python package \texttt{dbdreader}
(\url{https://github.com/smerckel/dbdreader}) is the de facto reader for
Slocum ocean glider binary data files: \texttt{.dbd},
\texttt{.sbd}, \texttt{.mbd} (engineering) and \texttt{.ebd}, \texttt{.tbd},
\texttt{.nbd} (science), plus their LZ4-compressed variants.  It is widely
used in operational glider data pipelines.

This document audits \texttt{dbdreader}'s implementation — the C
extension, the Python wrapper, and the LZ4 layer — identifies concrete
bugs and design problems, and describes a Julia replacement,
\texttt{JLDBDReader.jl}, validated byte-for-byte against \texttt{dbdreader}'s
output on real glider files.

\section{Architectural overview of \texttt{dbdreader}}

\texttt{dbdreader} is structured in four tiers:
\begin{enumerate}[leftmargin=*,topsep=1pt,itemsep=0pt]
  \item A C extension (\texttt{extension/dbdreader.c}, $\sim$770 lines)
        implementing the binary reader.
  \item A Python wrapper (\texttt{dbdreader/dbdreader.py}, $\sim$1900 lines)
        handling header parsing, cache files, parameter lookup,
        \texttt{MultiDBD} aggregation, and (importantly) a now-unused
        Python implementation of the binary reader, mirroring the C one.
  \item An LZ4 layer (\texttt{dbdreader/decompress.py}) using \texttt{python-lz4}.
  \item A scripts module providing CLI entry points.
\end{enumerate}

\section{Concrete defects}

\paragraph{1. \texttt{exit(1)} in the C extension.}
On a memory allocation failure or unsupported byte size, the C code calls
\texttt{exit(1)}.  This terminates the entire host Python interpreter
rather than returning an error to the caller.  Any pipeline embedding
\texttt{dbdreader} (e.g., a real-time data acquisition service) has no way
to handle the failure gracefully.

\paragraph{2. The \texttt{1e9} NaN sentinel.}
Missing values are encoded as \texttt{1e9} and detected with
\texttt{numpy.isclose(value, 1e9)}.  Two problems:
\begin{itemize}[leftmargin=*,topsep=1pt,itemsep=0pt]
  \item A real sensor reading numerically close to $10^9$ would be
        silently converted to NaN.
  \item Downstream code receives finite values that look real and only
        becomes NaN after a conditional check, breaking pipelines that
        assume IEEE NaN semantics.
\end{itemize}
The fix is to emit IEEE \texttt{NaN} directly, which is what
\texttt{JLDBDReader.jl} does.

\paragraph{3. NMEA validation gap: \texttt{minutes} $\geq$ 60.}
\texttt{dbdreader}'s NMEA-to-decimal conversion checks only that the
absolute value is below the degree limit ($90$ for latitude, $180$ for
longitude).  A value like \texttt{5360.0} (53\textdegree{}60\textquotesingle) passes the
check and converts to \texttt{54.0\textdegree{}} silently.  This happens in
practice when the glider's GPS reports malformed sentences.
\texttt{JLDBDReader.jl} adds a minutes-less-than-60 check.

\paragraph{4. \texttt{setlocale} mutation at module import.}
\texttt{dbdreader.py} calls \texttt{locale.setlocale(locale.LC\_ALL, "C")}
at top level.  Any code that runs \texttt{import dbdreader} after
configuring the locale (for \texttt{strftime}, currency, decimal
separators, etc.) is silently affected.

\paragraph{5. \texttt{mkdir} side-effect at module import.}
The last line of \texttt{dbdreader.py} instantiates a \texttt{DBDCache} object,
which creates a directory on disk.  Merely importing the package writes to
the user's filesystem.  In headless or read-only environments this can
fail or pollute paths.

\paragraph{6. Unreachable code with latent bugs.}
\texttt{dbdreader.py} contains a Python implementation of the binary
reader (\texttt{\_get\_by\_read\_per\_byte} and
\texttt{\_get\_by\_read\_per\_chunk}, $\sim$200 lines) that is no longer
called.  It byte-swaps unconditionally, never invoking the
endianness-detection logic; were it ever to be re-enabled, it would
silently corrupt data on most modern (little-endian) hosts.

\paragraph{7. 32-bit file offsets in C.}
The extension uses \texttt{int} for \texttt{ftell}-returned offsets
(\texttt{int fp\_end = ftell(\dots)}).  On 64-bit systems, files larger
than 2\,GiB are silently truncated.  Glider files rarely exceed this, but
\texttt{MultiDBD} can construct synthetic in-memory streams that do.

\paragraph{8. Thread-unsafe \texttt{static} variables.}
\texttt{read\_state\_bytes} in the C extension uses C \texttt{static}
locals for bit-shift constants.  Multiple Python threads calling into the
extension share these.  In current usage this happens to be safe because
the values never change, but the pattern is fragile and undocumented.

\paragraph{9. Stale file pointers in the Python class.}
\texttt{DBD.\_\_init\_\_} opens the file and stores
\texttt{self.fp}.  The handle is held open for the lifetime of the object,
even when no reads are happening, and is reused for multiple \texttt{seek}
operations across method calls — a fragile pattern that can leak handles
on exception paths and that interacts poorly with multiprocessing.

\paragraph{10. The cache-file-write text/binary mismatch.}
\texttt{DBDHeader.write\_cache\_file} opens the file in text mode but
writes bytes containing \texttt{\textbackslash{}n} line terminators.  On
Windows this produces \texttt{\textbackslash{}r\textbackslash{}n}, which
the reader then doesn't strip correctly.

\section{The \texttt{JLDBDReader.jl} reimplementation}

The Julia rewrite addresses every item above:

\begin{itemize}[leftmargin=*,topsep=1pt,itemsep=0pt]
  \item No C, hence no \texttt{exit(1)} and no compiler dependency.
  \item IEEE NaN throughout; no sentinel values.
  \item Strict NMEA validation including \texttt{minutes < 60}.
  \item No global state mutation, no import-time side effects.
  \item Fresh \texttt{open}/\texttt{close} per read; thread-safe by construction
        (immutable \texttt{DBDFile} struct, no shared mutable state).
  \item 64-bit offsets natively.
  \item Pure-Julia LZ4 block decoder, removing the \texttt{python-lz4} dependency.
  \item Built-in linear and heading-aware interpolation, removing \texttt{scipy}.
\end{itemize}

\subsection{Cycle layout (validated)}

Empirically verified by independent inspection of real \texttt{electa-2024-202-1-0.sbd}:

\begin{lstlisting}
17-byte known-cycle preamble:
  's' (0x73)  -- start tag
  1 byte int8 -- diagnostic
  uint16      -- 0x1234 endianness marker
  float32     -- 123.456
  float64     -- 123456789.12345
  'd' (0x64)  -- end tag

Each data cycle:
  state_bytes_per_cycle bytes (2 bits/sensor, MSB-first per byte)
  chunk_size bytes (sum of bytesizes for UPDATED sensors, in cycle order)
  1 separator byte                              <-- easily missed
\end{lstlisting}

\subsection{Validation against \texttt{dbdreader}}

A Python twin of the Julia algorithm was written and run against real
glider files from two independent gliders and four file types each:

\begin{itemize}[leftmargin=*,topsep=1pt,itemsep=0pt]
  \item Glider \texttt{electa}, deployment 2024-07-21:
        \texttt{02010000.\{dbd,sbd,tbd\}}.
  \item Glider \texttt{sylvia}, deployment 2024-09-30:
        \texttt{02390000.\{DBD,SBD,MBD,TBD\}}.
\end{itemize}

For each \texttt{(file, parameter)} combination the SHA-256 fingerprint of
the result's \texttt{float64} byte representation was computed and compared
against the output of \texttt{dbdreader} 0.5.x configured with the matching
cache files.

\paragraph{Result.}
All 34 \texttt{(file, parameter)} combinations tested produced
\emph{identical} float64 arrays — same length, same min/max, same SHA-256
— covering parameters such as \texttt{m\_depth}, \texttt{m\_heading},
\texttt{m\_pitch}, \texttt{m\_roll}, \texttt{m\_battery},
\texttt{m\_gps\_lat}, \texttt{m\_gps\_lon}, \texttt{m\_present\_time},
\texttt{sci\_water\_temp}, \texttt{sci\_water\_cond},
\texttt{sci\_water\_pressure}, and \texttt{sci\_m\_present\_time}.  The
two gliders span different sensor list CRCs, different
\texttt{sensors\_per\_cycle} values (102 / 1729 / 134 / 10 for sylvia vs.
64 / 1696 / 134 / 14 for electa), and different mission configurations
— a reasonable cross-section.

\subsection{Bugs found during \texttt{JLDBDReader.jl} development}

Self-disclosure: an earlier draft (released as \texttt{SlocumGliders.jl})
had two bugs caught during real-data validation:

\begin{enumerate}[leftmargin=*,topsep=1pt,itemsep=0pt]
  \item \texttt{chunk\_offset} was accumulated only over requested sensors;
    must be over all UPDATED sensors so that the byte offset of each value
    within the chunk is correct.
  \item \texttt{SensorInfo.index} was set to \texttt{full\_idx} (position in
    the master sensor namespace) rather than \texttt{active\_pos} (position
    in the cycle layout).  These usually coincide because cache files list
    sensors in full-index order, but conflating them is brittle.
\end{enumerate}

Both are fixed in the current package, and the integration test
(\texttt{test/runtests.jl}) compares against the reference fingerprints to
prevent regressions.

\subsection{Limitations and honest disclosures}

\begin{itemize}[leftmargin=*,topsep=1pt,itemsep=0pt]
  \item Without a matching \texttt{.cac} cache file, no DBD-family file can
        be decoded.  This is a fundamental property of the format
        (\texttt{sensor\_list\_factored: 1}), not a package limitation.
  \item The \texttt{*LIST.DAT} mission-config files (\texttt{SBDLIST.DAT}
        etc.) are not substitutes for cache files: they record
        \emph{recording-interval} metadata, not byte sizes.
  \item The MBD and EBD files for glider \texttt{electa} in the original
        fixture set lacked matching cache files (\texttt{9cde8b23} and
        \texttt{085b9e63}) and could not be byte-validated.  Sylvia's
        full file set (DBD/SBD/MBD/TBD) has been validated.
  \item Encoding versions other than \texttt{5} have not been tested
        (the package issues a warning).
  \item The package targets Julia $\geq$ 1.10.
\end{itemize}

\subsection{What the package does \emph{not} do}

\texttt{JLDBDReader.jl} reads the raw binary; it does not provide:
\begin{itemize}[leftmargin=*,topsep=1pt,itemsep=0pt]
  \item Higher-level products: calibrated CTD, dive segmentation, glider
        flight model.  These belong in a downstream package.
  \item Cache-file generation from inline sensor lists.  (Easy to add, but
        out of scope for the initial port.)
  \item Writing/modifying glider files.  Read-only by design.
\end{itemize}

\end{document}
